% =====================================================================
% ejercicio01.tex
% Responsabilidad única: derivar las ecuaciones normales de la recta de
% mínimos cuadrados, ajustar los datos didácticos, verificar propiedades
% de los residuos y estudiar el efecto de una observación influyente.
% Cifras verificadas con verificacion/teoria.py.
% =====================================================================
\subsection{Ejercicio 1: Regresión lineal simple, ajuste e influencia de una observación}

\paragraph{La situación.} Una empresa quiere saber cuánta energía consume
un equipo según las horas que pasa encendido. Para cada uno de $n=5$ días
anotamos las horas de funcionamiento $x_i$ y el consumo $y_i$ en kWh:

\begin{table}[H]
\centering
\begin{tabular}{c|ccccc}
\toprule
Día $i$ & 1 & 2 & 3 & 4 & 5 \\
\midrule
Horas $x_i$ & 1 & 2 & 3 & 4 & 5 \\
Consumo $y_i$ (kWh) & 2 & 4 & 5 & 4 & 5 \\
\bottomrule
\end{tabular}
\caption{Datos del Ejercicio 1.}
\label{tab:ej1-datos}
\end{table}

Queremos resumir la relación con una recta
$\hat y_i=\beta_0+\beta_1x_i$, donde $\beta_0$ (el \emph{intercepto}) es la
altura de la recta cuando $x=0$ y $\beta_1$ (la \emph{pendiente}) es cuánto
sube el consumo predicho por cada hora adicional. Ninguna recta pasa por
los cinco puntos, así que cada predicción tendrá un error
$y_i-\hat y_i$. El criterio de \emph{mínimos cuadrados} elige la recta que
hace más pequeña la suma de esos errores elevados al cuadrado,
\[
    S(\beta_0,\beta_1)=\sum_{i=1}^{n}\bigl(y_i-\beta_0-\beta_1x_i\bigr)^2 .
\]
Elevar al cuadrado cumple dos funciones: evita que los errores positivos y
negativos se cancelen entre sí y castiga mucho más un error grande que
varios pequeños (un error de $4$ cuenta $16$, mientras que dos errores de
$2$ cuentan $4+4=8$).

\begin{enumerate}

\item \textbf{Ecuaciones normales y estimadores.}
$S$ es una función de dos variables, $\beta_0$ y $\beta_1$, que tiene forma
de tazón. En el fondo del tazón la superficie es plana en todas las
direcciones, así que ahí las dos derivadas parciales valen cero. Para
derivar usamos la regla de la cadena: la derivada de $(u)^2$ es $2u$
multiplicado por la derivada de $u$.

\emph{Paso 1: derivada respecto a $\beta_0$.} Dentro del paréntesis,
$u_i=y_i-\beta_0-\beta_1x_i$, cuya derivada respecto a $\beta_0$ es $-1$,
luego
\[
    \frac{\partial S}{\partial \beta_0}
    = \sum_i 2\,(y_i-\beta_0-\beta_1x_i)\cdot(-1)
    = -2\sum_i (y_i-\beta_0-\beta_1x_i) .
\]

\emph{Paso 2: derivada respecto a $\beta_1$.} Ahora la derivada de $u_i$
respecto a $\beta_1$ es $-x_i$, luego
\[
    \frac{\partial S}{\partial \beta_1}
    = \sum_i 2\,(y_i-\beta_0-\beta_1x_i)\cdot(-x_i)
    = -2\sum_i x_i(y_i-\beta_0-\beta_1x_i) .
\]

\emph{Paso 3: igualar a cero y ordenar.} Al igualar ambas a cero podemos
dividir entre $-2$ y repartir la suma término a término (recordando que
$\sum_i\beta_0=n\beta_0$, pues $\beta_0$ se suma $n$ veces):
\begin{align*}
    \sum_i y_i - n\beta_0 - \beta_1\sum_i x_i = 0
    &\quad\Longrightarrow\quad
    n\beta_0 + \beta_1\sum_i x_i = \sum_i y_i, \\
    \sum_i x_iy_i - \beta_0\sum_i x_i - \beta_1\sum_i x_i^2 = 0
    &\quad\Longrightarrow\quad
    \beta_0\sum_i x_i + \beta_1\sum_i x_i^2 = \sum_i x_iy_i .
\end{align*}
Estas dos igualdades son las \emph{ecuaciones normales}: un sistema de dos
ecuaciones con dos incógnitas.

\emph{Paso 4: despejar $\beta_0$.} Dividimos la primera ecuación entre $n$
y usamos que $\frac1n\sum_i x_i=\bar x$ y $\frac1n\sum_i y_i=\bar y$:
\[
    \beta_0 + \beta_1\bar x = \bar y
    \quad\Longrightarrow\quad
    \beta_0=\bar y-\beta_1\bar x .
\]
Esto dice algo útil por sí solo: la recta ajustada siempre pasa por el
punto de promedios $(\bar x,\bar y)$, porque al reemplazar $x=\bar x$
queda $\hat y=\bar y-\beta_1\bar x+\beta_1\bar x=\bar y$.

\emph{Paso 5: despejar $\beta_1$.} Sustituimos $\beta_0=\bar y-\beta_1\bar x$
en la segunda ecuación y usamos $\sum_i x_i=n\bar x$:
\begin{align*}
    (\bar y-\beta_1\bar x)\,n\bar x + \beta_1\sum_i x_i^2 &= \sum_i x_iy_i \\
    n\bar x\bar y - \beta_1 n\bar x^2 + \beta_1\sum_i x_i^2 &= \sum_i x_iy_i \\
    \beta_1\Bigl(\sum_i x_i^2 - n\bar x^2\Bigr) &= \sum_i x_iy_i - n\bar x\bar y .
\end{align*}

\emph{Paso 6: escribirlo con desviaciones.} Las dos expresiones entre
paréntesis se pueden escribir de forma más intuitiva. Expandiendo el
cuadrado y usando de nuevo $\sum_i x_i=n\bar x$,
\[
    \sum_i (x_i-\bar x)^2
    = \sum_i x_i^2 - 2\bar x\sum_i x_i + n\bar x^2
    = \sum_i x_i^2 - 2n\bar x^2 + n\bar x^2
    = \sum_i x_i^2 - n\bar x^2 ,
\]
y de la misma forma
\[
    \sum_i (x_i-\bar x)(y_i-\bar y)
    = \sum_i x_iy_i - \bar y\sum_i x_i - \bar x\sum_i y_i + n\bar x\bar y
    = \sum_i x_iy_i - n\bar x\bar y .
\]
Llamamos $S_{xx}=\sum_i(x_i-\bar x)^2$ (qué tanto varían las horas
alrededor de su promedio) y $S_{xy}=\sum_i(x_i-\bar x)(y_i-\bar y)$ (qué
tanto varían juntas horas y consumo: es positivo si cuando $x$ está por
encima de su promedio $y$ también suele estarlo). Con esto llegamos a
\[
    \boxed{\;
    \hat\beta_1 = \frac{\sum_i (x_i-\bar x)(y_i-\bar y)}{\sum_i (x_i-\bar x)^2}
    = \frac{S_{xy}}{S_{xx}},
    \qquad
    \hat\beta_0 = \bar y - \hat\beta_1\bar x .\;}
\]

\emph{Condición para que la pendiente sea única.} La fórmula de
$\hat\beta_1$ solo tiene sentido si $S_{xx}>0$, es decir, si los valores de
$x$ \textbf{no son todos iguales}. Si todos los días el equipo hubiera
funcionado las mismas horas, todos los puntos estarían en una misma línea
vertical y no habría forma de saber cómo cambia el consumo al cambiar las
horas: cualquier recta que pase por $(\bar x,\bar y)$ daría el mismo error.

\emph{Comprobación de que es un mínimo.} Que las derivadas valgan cero
también ocurre en un máximo o en un punto de silla, así que revisamos las
segundas derivadas, que se organizan en la matriz hessiana:
\[
    \frac{\partial^2 S}{\partial \beta_0^2}=2n,\qquad
    \frac{\partial^2 S}{\partial \beta_0\partial\beta_1}=2\sum_i x_i,\qquad
    \frac{\partial^2 S}{\partial \beta_1^2}=2\sum_i x_i^2
    \quad\Longrightarrow\quad
    H = 2\begin{bmatrix} n & \sum_i x_i\\ \sum_i x_i & \sum_i x_i^2\end{bmatrix}.
\]
Su primer elemento $2n$ es positivo y su determinante es
$4\bigl(n\sum_i x_i^2-(\sum_i x_i)^2\bigr)=4n\bigl(\sum_i x_i^2-n\bar x^2\bigr)=4nS_{xx}$,
también positivo cuando $S_{xx}>0$. Una matriz $2\times2$ con esas dos
condiciones es \emph{definida positiva}, lo que en palabras quiere decir
que $S$ se curva hacia arriba en todas las direcciones (el tazón), así que
el punto que encontramos es el mínimo global.

\item \textbf{Coeficientes, predicciones y residuos.}
\emph{Paso 1: promedios.} $\bar x=(1+2+3+4+5)/5=15/5=3$ horas y
$\bar y=(2+4+5+4+5)/5=20/5=4$ kWh.

\emph{Paso 2: desviaciones respecto al promedio.}
$x_i-\bar x=(-2,-1,0,1,2)$ y $y_i-\bar y=(-2,0,1,0,1)$.

\emph{Paso 3: $S_{xy}$ y $S_{xx}$.} Multiplicamos las desviaciones día a
día y sumamos:
\[
    S_{xy} = (-2)(-2)+(-1)(0)+(0)(1)+(1)(0)+(2)(1) = 4+0+0+0+2 = 6,
    \qquad
    S_{xx} = 4+1+0+1+4 = 10 .
\]

\emph{Paso 4: coeficientes.}
\[
    \hat\beta_1 = \frac{6}{10} = 0.6, \qquad
    \hat\beta_0 = 4 - 0.6\cdot 3 = 2.2 .
\]
Como control, estos valores cumplen las ecuaciones normales del inciso a
con $\sum_i x_i=15$, $\sum_i x_i^2=55$, $\sum_i y_i=20$ y
$\sum_i x_iy_i=66$: $5(2.2)+15(0.6)=20$ y $15(2.2)+55(0.6)=66$.

\emph{Paso 5: predicciones y residuos.} Con la recta
$\hat y=2.2+0.6x$ calculamos la predicción de cada día y el residuo
$e_i=y_i-\hat y_i$, que es lo que la recta no alcanzó a explicar
(Tabla~\ref{tab:ej1-residuos}).

\begin{table}[H]
\centering
\begin{tabular}{crrrrr}
\toprule
$i$ & $x_i$ & $y_i$ & $\hat y_i = 2.2+0.6x_i$ & $e_i=y_i-\hat y_i$ & $e_i^2$ \\
\midrule
1 & 1 & 2 & 2.8 & $-0.8$ & 0.64 \\
2 & 2 & 4 & 3.4 & $0.6$  & 0.36 \\
3 & 3 & 5 & 4.0 & $1.0$  & 1.00 \\
4 & 4 & 4 & 4.6 & $-0.6$ & 0.36 \\
5 & 5 & 5 & 5.2 & $-0.2$ & 0.04 \\
\midrule
$\sum$ & & & & $0.0$ & 2.40 \\
\bottomrule
\end{tabular}
\caption{Predicciones y residuos del ajuste $\hat y = 2.2 + 0.6x$.}
\label{tab:ej1-residuos}
\end{table}

\emph{Paso 6: los residuos suman cero.} En efecto
$\sum_ie_i=-0.8+0.6+1.0-0.6-0.2=0$, y al revisarlo vimos que no es una
coincidencia de estos datos sino la primera ecuación normal disfrazada.
Esa ecuación salió de $\partial S/\partial\beta_0=0$, que es
\[
    \sum_i(y_i-\hat\beta_0-\hat\beta_1x_i)=\sum_i(y_i-\hat y_i)=\sum_ie_i=0 .
\]
En palabras: como el intercepto es libre de moverse, el ajuste lo sube o
lo baja hasta que los errores por encima y por debajo de la recta se
compensan exactamente. Por eso la propiedad depende de incluir el
intercepto: si ajustáramos una recta obligada a pasar por el origen,
$\hat y=\beta_1x$, no existiría $\beta_0$, desaparecería esa ecuación y
solo quedaría $\sum_ix_ie_i=0$, de modo que en general los residuos ya no
sumarían cero.

\item \textbf{MSE, $R^2$ y comparación con $\bar y$.}
\emph{Paso 1: error cuadrático medio.} Es el promedio de los errores al
cuadrado de la Tabla~\ref{tab:ej1-residuos}:
\[
    \mathrm{MSE} = \frac{1}{n}\sum_i e_i^2 = \frac{2.4}{5} = 0.48 \ \text{kWh}^2 .
\]

\emph{Paso 2: el punto de comparación.} El modelo más ingenuo posible
ignora las horas y predice siempre el consumo promedio, $\hat y=\bar y=4$.
Sus errores son las desviaciones $y_i-\bar y=(-2,0,1,0,1)$, cuya suma de
cuadrados es la variabilidad total
\[
    \mathrm{SST}=\sum_i (y_i-\bar y)^2 = 4+0+1+0+1 = 6,
    \qquad
    \mathrm{MSE}_{\bar y}=\frac{6}{5}=1.2\ \text{kWh}^2 .
\]

\emph{Paso 3: coeficiente de determinación.} $R^2$ compara cuánto error
le queda a la recta frente al que tenía el promedio:
\[
    R^2 = 1 - \frac{\mathrm{SSE}}{\mathrm{SST}} = 1 - \frac{2.4}{6} = 1-0.4 = 0.6 .
\]
La recta reduce el error cuadrático medio de $1.2$ a $0.48$, es decir, en
un 60\,\%, y eso es justamente lo que mide $R^2$: el 60\,\% de la
variabilidad del consumo alrededor de su promedio queda explicada por las
horas de funcionamiento, y el 40\,\% restante se debe a otras cosas que el
modelo no ve.

\emph{Paso 4: interpretación con unidades.} La pendiente tiene unidades de
kWh por hora, porque divide un cambio de consumo entre un cambio de horas.
Así, $\hat\beta_1=0.6$ significa que, en estos datos, \textbf{cada hora
adicional de funcionamiento se asocia con un aumento promedio de $0.6$ kWh}
en el consumo. El intercepto, $2.2$ kWh, sería el consumo predicho con el
equipo encendido cero horas, pero ese valor está fuera del rango que
observamos ($1\le x\le5$), así que preferimos no leerlo como un consumo
real en reposo: es solo la altura a la que la recta corta el eje.

\item \textbf{Observación influyente ($y_5 = 12$).}
Ahora cambiamos un solo dato: el consumo del día 5 pasa de $5$ a $12$ kWh,
y repetimos el ajuste.

\emph{Paso 1: nuevo promedio de $y$.} $\bar y'=(2+4+5+4+12)/5=27/5=5.4$.
Las $x$ no cambiaron, así que $\bar x=3$ y $S_{xx}=10$ siguen igual.

\emph{Paso 2: nuevo $S_{xy}$.} Aquí usamos un atajo: como las desviaciones
de $x$ suman cero, $\sum_i(x_i-\bar x)=0$, el término con $\bar y$
desaparece,
\[
    S_{xy}=\sum_i(x_i-\bar x)(y_i-\bar y)
    =\sum_i(x_i-\bar x)y_i-\bar y\underbrace{\sum_i(x_i-\bar x)}_{=0}
    =\sum_i(x_i-\bar x)y_i ,
\]
y basta multiplicar cada desviación de $x$ por el nuevo $y$:
\[
    S_{xy}' = (-2)(2)+(-1)(4)+(0)(5)+(1)(4)+(2)(12) = -4-4+0+4+24 = 20 .
\]

\emph{Paso 3: nuevos coeficientes y predicciones en $x=6$.}
\[
    \hat\beta_1' = \frac{20}{10} = 2, \qquad
    \hat\beta_0' = 5.4 - 2\cdot 3 = -0.6,
\]
\[
    \hat y'(6) = -0.6 + 2\cdot 6 = 11.4
    \quad\text{frente a}\quad
    \hat y(6) = 2.2 + 0.6\cdot 6 = 5.8 .
\]
Nos sorprendió que una sola observación multiplicara la pendiente por más
de tres (de $0.6$ a $2$) y casi duplicara la predicción en $x=6$ (figura~\ref{fig:ej1-rectas}).

\emph{Paso 4: por qué la pérdida cuadrática reacciona tanto.} Con la recta
original, el punto $(5,12)$ tendría residuo $12-5.2=6.8$ y aportaría
$6.8^2=46.24$ al SSE, mucho más que los otros cuatro residuos juntos, que
suman $2.36$. Como el criterio eleva al cuadrado, ese único error domina la
suma, y al ajuste le conviene inclinar toda la recta para reducirlo aunque
empeore a los demás. Tras el reajuste, el residuo de ese punto baja de
$6.8$ a $2.6$ (su cuadrado de $46.24$ a $6.76$), mientras que el punto
$(4,4)$ pasa de $-0.6$ a $-3.4$, y aun así el SSE sobre los datos
modificados cae de $48.6$ a $19.2$ (figura~\ref{fig:ej1-cuadrados}).

\emph{Paso 5: por qué importa dónde está el punto.} El punto modificado
está en el extremo derecho del rango de $x$. Como vimos en el paso 2,
$\hat\beta_1=\sum_i(x_i-\bar x)y_i/S_{xx}$, así que si solo cambia $y_5$
en $\Delta y_5=12-5=7$, la pendiente cambia en
\[
    \Delta\hat\beta_1 = \frac{(x_5-\bar x)\,\Delta y_5}{S_{xx}}
    = \frac{2\cdot7}{10}=1.4 ,
\]
que es exactamente $2-0.6$. El peso de cada punto sobre la pendiente es
proporcional a su distancia horizontal al promedio, $x_i-\bar x$, como
una palanca: entre más lejos del centro, más fuerza hace. Esa idea se mide
con el \emph{apalancamiento}
$h_i=\frac1n+\frac{(x_i-\bar x)^2}{S_{xx}}$, que aquí vale
$h_5=\frac15+\frac{4}{10}=0.6$ para el día 5 frente a $h_3=\frac15+0=0.2$
para el día 3. Por eso el mismo cambio en $y_3$ (con $x_3=\bar x$) no
habría movido la pendiente en absoluto, solo habría subido la recta.

\emph{Paso 6: el problema de predecir en $x=20$.} Predecir para $20$ horas
es una \textbf{extrapolación}: cuatro veces el máximo que observamos. Nada
en los datos garantiza que la relación siga siendo una recta fuera de
$[1,5]$, y además allí los dos modelos discrepan enormemente
($\hat y(20)=2.2+0.6\cdot20=14.2$ frente a
$\hat y'(20)=-0.6+2\cdot20=39.4$). La razón es que cada recta pasa por su
propio punto de promedios, de modo que la diferencia entre ambas es
\[
    \hat y'(x)-\hat y(x)
    =\underbrace{(\bar y'-\bar y)}_{=1.4}
    +\underbrace{(\hat\beta_1'-\hat\beta_1)}_{=1.4}\,(x-\bar x)
    =1.4+1.4(x-3) ,
\]
y la parte que viene de la pendiente crece en línea recta con la distancia
$|x-\bar x|$: la brecha es $5.6$ kWh en $x=6$ y $25.2$ kWh en $x=20$. Una
pequeña duda sobre la pendiente se vuelve una duda enorme lejos de los
datos.

\end{enumerate}

\begin{figure}[H]
\centering
\begin{tikzpicture}
\begin{axis}[
    width=0.72\textwidth, height=6.2cm,
    xmin=0.5, xmax=6.5, ymin=0, ymax=13,
    xlabel={$x$ (horas de funcionamiento)}, ylabel={$y$ (kWh)},
    legend pos=north west, legend cell align=left,
    legend style={font=\small}, tick label style={font=\small},
]
\addplot[only marks, mark=*, color=classblue, mark size=2.4pt]
    coordinates {(1,2) (2,4) (3,5) (4,4) (5,5)};
\addlegendentry{datos originales}
\addplot[only marks, mark=square*, color=classred, mark size=2.6pt]
    coordinates {(5,12)};
\addlegendentry{$y_5$ modificado ($=12$)}
\addplot[domain=0.5:6.5, thick, color=classblue] {2.2 + 0.6*x};
\addlegendentry{$\hat y = 2.2 + 0.6x$}
\addplot[domain=0.5:6.5, thick, dashed, color=classred] {-0.6 + 2*x};
\addlegendentry{$\hat y' = -0.6 + 2x$}
\addplot[only marks, mark=o, color=black, mark size=2.4pt]
    coordinates {(6,5.8) (6,11.4)};
\end{axis}
\end{tikzpicture}
\caption{Recta de mínimos cuadrados con los datos originales (azul) y
tras cambiar solo $y_5$ a $12$ (rojo); los círculos vacíos marcan las
predicciones en $x=6$ ($5.8$ y $11.4$).}
\label{fig:ej1-rectas}
\end{figure}

\begin{figure}[H]
\centering
\begin{tikzpicture}
\begin{axis}[
    width=0.8\textwidth, height=5.8cm,
    ybar, bar width=11pt, enlarge x limits=0.12,
    ymin=0, ymax=60, area legend,
    xtick={1,2,3,4,5},
    xticklabels={{$(1,2)$},{$(2,4)$},{$(3,5)$},{$(4,4)$},{$(5,12)$}},
    xlabel={observación $(x_i,y_i)$ con $y_5=12$},
    ylabel={$e_i^2$ (kWh$^2$)},
    nodes near coords, every node near coord/.append style={font=\scriptsize,
        rotate=90, anchor=west, /pgf/number format/fixed,
        /pgf/number format/fixed zerofill, /pgf/number format/precision=2},
    legend pos=north west, legend cell align=left,
    legend style={font=\small}, tick label style={font=\small},
]
\addplot[fill=classblue!70, draw=classblue]
    coordinates {(1,0.64) (2,0.36) (3,1.00) (4,0.36) (5,46.24)};
\addlegendentry{recta original $2.2+0.6x$ (SSE $=48.6$)}
\addplot[fill=classred!70, draw=classred]
    coordinates {(1,0.36) (2,0.36) (3,0.16) (4,11.56) (5,6.76)};
\addlegendentry{recta reajustada $-0.6+2x$ (SSE $=19.2$)}
\end{axis}
\end{tikzpicture}
\caption{Contribución de cada observación al SSE de los datos modificados
($y_5=12$), evaluada con la recta original y con la reajustada; bajo la
recta original el punto $(5,12)$ concentra casi todo el error, de modo
que mínimos cuadrados acepta empeorar los residuos de los demás puntos
(en especial el de $(4,4)$) a cambio de reducir ese único cuadrado de
$46.24$ a $6.76$.}
\label{fig:ej1-cuadrados}
\end{figure}
